\section*{Appendix S2. M1 Isolated-Band Calculation}
\label{supp:m1-isolated-band}
\renewcommand{\theequation}{S2.\arabic{equation}}
\setcounter{equation}{0}
\renewcommand{\thetable}{S2.\arabic{table}}
\setcounter{table}{0}

\subsection*{S2.1 Role and frozen scope}

M1 is the planning identifier for the prospectively frozen, controlled
one-dimensional lowest-band calculation used in the main paper. Its retained
calculation identity is
\begin{center}
\path{icmsep2026.paper1.periodic1d.isolated-band.v1}.
\end{center}
The calculation is synthetic numerical-verification evidence. It evaluates a
finite parent representation, extracts the lowest scalar band, performs a
complete reciprocal-to-hopping transform, constructs finite-range reductions,
compares matched routes, and evaluates disjoint withheld points. The retained
schema contains no band-gap result and does not by itself establish that this
lowest band is isolated; a sampled-gap or stronger isolation argument is a
separate prerequisite channel. The historical identity is not reinterpreted as
such evidence. The calculation contains no material parameters or external
electronic-structure execution.

The dimensionless parent is
\begin{equation}
 \widehat H(k)=-\frac{\mathrm d^2}{\mathrm d x^2}+0.5\cos x,
 \qquad a=2\pi,\quad G=1,\quad E_G=1.
 \label{eq:supp-m1-parent}
\end{equation}
Reduced momentum $k/G$ lies in $[-1/2,1/2]$. The energy zero is fixed by the
model; no post-calculation energy shift is fitted. Table~\ref{tab:supp-m1-controls}
collects the controls frozen before confirmatory execution.

\begin{table}[ht]
\centering
\caption{Prospectively frozen M1 controls.}
\label{tab:supp-m1-controls}
\begin{tabular}{@{}ll@{}}
\toprule
Control & Frozen value \\
\midrule
Plane-wave cutoffs & $3,5,7,9,11$ \\
Finite represented reference cutoff & $15$ \\
Production plane-wave cutoff & $11$ \\
Finite-difference grids & $31,63,127,255$ points \\
Parent momenta & $-1/2,-1/4,0,1/4,1/2$ \\
Compared parent bands & lowest three ordered bands \\
Training mesh & 64-point centered half-open mesh \\
Hopping ranges & $0,1,2,3,4,6,8$ \\
Withheld mesh & 257-point staggered uniform mesh \\
Numerical-consistency tolerance & $10^{-12}$ \\
Coordinate tolerance & $10^{-14}$ \\
\bottomrule
\end{tabular}
\end{table}

The historical Mathieu values, common-low-mode comparison, weak-potential
sweep, and stress attacks are outside M1. Multiband alignment, nonidentity gauge
families, Wannier localization, topology, two-dimensional shells, material
validation, and uncertainty quantification are also outside its frozen scope.

\subsection*{S2.2 Parent-representation diagnostics}

At each frozen momentum, the lowest three eigenvalues were compared with the
cutoff-15 plane-wave representation. This cutoff is the declared finite
reference for M1, not an exact continuum spectrum. The plane-wave observations
are listed in Table~\ref{tab:supp-m1-plane-wave}.

\begin{table}[ht]
\centering
\caption{Maximum absolute plane-wave spectral discrepancy relative to the
cutoff-15 represented reference.}
\label{tab:supp-m1-plane-wave}
\begin{tabular}{@{}rr@{}}
\toprule
Cutoff $P$ & Maximum discrepancy ($E_G$) \\
\midrule
3  & $1.2107868154753731\times10^{-5}$ \\
5  & $5.311306949806749\times10^{-13}$ \\
7  & $6.483702463810914\times10^{-14}$ \\
9  & $8.038014698286133\times10^{-14}$ \\
11 & $8.448797217397441\times10^{-14}$ \\
\bottomrule
\end{tabular}
\end{table}

The cutoff-7 through cutoff-11 discrepancies fluctuate at the represented
binary64 scale, so this sequence is not interpreted as a monotone continuum
error estimate. The independently assembled centered finite-difference channel
produced Table~\ref{tab:supp-m1-fd}.

\begin{table}[ht]
\centering
\caption{Maximum absolute finite-difference spectral discrepancy relative to
the same cutoff-15 represented reference.}
\label{tab:supp-m1-fd}
\begin{tabular}{@{}rr@{}}
\toprule
Grid points & Maximum discrepancy ($E_G$) \\
\midrule
31  & $1.75327419749558\times10^{-2}$ \\
63  & $4.2534045602447\times10^{-3}$ \\
127 & $1.0471636947908536\times10^{-3}$ \\
255 & $2.597716946719508\times10^{-4}$ \\
\bottomrule
\end{tabular}
\end{table}

The observed decrease is the finite-grid refinement diagnostic. M1 does not
convert it into a continuum convergence certificate.

\subsection*{S2.3 Training and withheld roles}

The 64-point training mesh supplies the complete discrete Fourier transform and
the equal-weight direct fits. For training extent $N=64$ and withheld extent
$M=257$, withheld coordinate $i$ is
\begin{equation}
 \frac{k_i}{G}=-\frac12+\frac{i+1/(N+1)}{M},
 \qquad i=0,\ldots,M-1.
 \label{eq:supp-m1-withheld}
\end{equation}
The retained verifier confirms that the training and withheld coordinates are
disjoint. Withheld values enter evaluation only; they do not alter coefficients,
ranges, tolerances, or figure selection.

The complete scalar hopping coefficients are
\begin{equation}
 t_R=\frac{1}{N}\sum_{j=0}^{N-1}
 e^{-2\pi iRk_j/G}E(k_j).
 \label{eq:supp-m1-transform}
\end{equation}
For the centered half-open mesh $k_j/G=-1/2+j/N$, the Fourier columns obey
\begin{equation}
 \frac{1}{N}\sum_{j=0}^{N-1}
 e^{2\pi i(S-R)k_j/G}
 =(-1)^{S-R}\,\delta_{R,S\ (\mathrm{mod}\ N)}.
 \label{eq:supp-m1-fourier-orthogonality}
\end{equation}
Thus distinct retained representatives modulo $N$ are orthogonal; the centered
origin contributes the factor $(-1)^{S-R}$. Their inverse transform reconstructs
the training samples with maximum absolute error
$4.166137705498387\times10^{-17}E_G$. The maximum modular Hermiticity defect is
$2.0907097393389345\times10^{-18}E_G$.

For each range $R_{\max}$, the mediated route truncates the complete
coefficients to $|R|\leq R_{\max}$. The direct route fits those same
representatives by equal-weight complex least squares on the same training
mesh. Tables~\ref{tab:supp-m1-ranges} and
\ref{tab:supp-m1-route-defects} report the retained range diagnostics.

\begin{table}[ht]
\centering
\small
\caption{M1 finite-range error and omitted-block diagnostics. All diagnostic
columns have energy unit $E_G$.}
\label{tab:supp-m1-ranges}
\begin{tabular}{@{}rrrr@{}}
\toprule
$R_{\max}$ & Training maximum & Withheld maximum & Omitted-block norm \\
\midrule
0 & $4.6235761235\times10^{-2}$ & $4.6235757038\times10^{-2}$ & $3.0304371874\times10^{-2}$ \\
1 & $3.4907261092\times10^{-3}$ & $3.4907249362\times10^{-3}$ & $2.1876792983\times10^{-3}$ \\
2 & $4.1809353421\times10^{-4}$ & $4.1809323067\times10^{-4}$ & $2.5574448521\times10^{-4}$ \\
3 & $6.0054872741\times10^{-5}$ & $6.0054797140\times10^{-5}$ & $3.6185634797\times10^{-5}$ \\
4 & $9.5136496848\times10^{-6}$ & $9.5136312847\times10^{-6}$ & $5.6738480443\times10^{-6}$ \\
6 & $2.8158443985\times10^{-7}$ & $2.8158339337\times10^{-7}$ & $1.6575039069\times10^{-7}$ \\
8 & $9.4653549651\times10^{-9}$ & $9.4652957590\times10^{-9}$ & $5.5272060152\times10^{-9}$ \\
\bottomrule
\end{tabular}
\end{table}

\begin{table}[ht]
\centering
\caption{M1 direct-versus-mediated coefficient defects.}
\label{tab:supp-m1-route-defects}
\begin{tabular}{@{}rr@{}}
\toprule
$R_{\max}$ & Coefficient defect ($E_G$) \\
\midrule
0 & $4.1633363423\times10^{-17}$ \\
1 & $3.4428416195\times10^{-17}$ \\
2 & $5.9106172771\times10^{-17}$ \\
3 & $7.3812041610\times10^{-17}$ \\
4 & $6.0865104451\times10^{-17}$ \\
6 & $4.7557874187\times10^{-17}$ \\
8 & $6.4715834978\times10^{-17}$ \\
\bottomrule
\end{tabular}
\end{table}

Across the frozen sequence, the largest direct-versus-mediated coefficient
defect is $7.38120416098117\times10^{-17}E_G$, and the largest sampled maximum
route defect is $1.564917933602745\times10^{-16}E_G$. The largest Parseval
residual is $6.938893903907228\times10^{-18}E_G^2$, and the largest imaginary
band-shape residual is $9.479926314907974\times10^{-16}E_G$. These are matched-
route and numerical-consistency diagnostics. M1 defines no preferred hopping
range or scientific acceptance threshold.

\subsection*{S2.4 Independent reconstruction and provenance}

The retained verifier strictly decodes the frozen input and result documents
and imports neither the producer entry point nor the maintained
\texttt{ksdft2effmass} calculation implementation. It independently rebuilds
parent matrices, spectra, Fourier coefficients, truncations, least-squares
fits, route defects, Hermiticity, Parseval quantities, and band-shape
diagnostics. It nevertheless shares NumPy, SciPy, eigensolver semantics,
model conventions, and the runtime environment with the producer; it is an
independent finite-protocol reconstruction, not an independent physical or
continuum oracle. The verification passed with the defects in
Table~\ref{tab:supp-m1-verification}.

\begin{table}[ht]
\centering
\caption{Independent retained-reconstruction defects.}
\label{tab:supp-m1-verification}
\begin{tabular}{@{}lr@{}}
\toprule
Reconstructed channel & Maximum absolute defect \\
\midrule
Spectral & $0$ \\
Hopping & $4.781392881215698\times10^{-17}$ \\
Diagnostic & $2.731148640577885\times10^{-14}$ \\
\bottomrule
\end{tabular}
\end{table}

All are below the frozen $10^{-12}$ verification tolerance. The result and
verification wire identities are
\begin{center}
\path{ksdft2effmass.periodic1d.isolated-band-calculation-result.v1},\\
\path{ksdft2effmass.periodic1d.isolated-band-calculation-verification.v1}.
\end{center}
The protocol-freeze record contains the hashes of \path{input.json} and
\path{protocol.md} recorded before confirmatory execution. The retained package
further binds the result, verification, scripts, software record, source
manifest, figure data, and figure through \path{SHA256SUMS}. These digests test
byte identity relative to the retained manifests; by themselves they do not
supply an external trusted timestamp or prove chronology. The source manifest
records the calculation snapshot and is not a claim that a later evolving
worktree has identical hashes. The calculation used local in-process Python
software only; it invoked no external electronic-structure calculator,
scheduler, remote resource, or material dataset.
